\subsection{APPLICATION: ORDER OF THE WEYL GROUP}

Lemma 1. Let X be a locally compact space countable at infinity, G a group acting continuously and properly on X, \(\mu\) a measure on X invariant under G, G' a subgroup of G, U and U' two open subsets of X with finite non-zero measure. Assume that the sU for \(s \in G\) (resp. the \(s'U'\) for \(s' \in G'\) ) are pairwise disjoint and that their union is of negligible complement. Then G' is of finite index in G and \((\mathrm{G} : \mathrm{G}') = \mu(\mathrm{U}') / \mu(\mathrm{U})\) .

Let \((s_{\lambda})_{\lambda \in \Lambda}\) be a family of representatives of the right cosets of \(G'\) in G. Let \(U_1\) be the union of the \(s_{\lambda}U\) . Then the \(s'U_1\) , for \(s' \in G'\) , are pairwise disjoint and have union \(M = \bigcup_{s \in G} sU\) . Let \(M' = \bigcup_{s' \in G'} s'U'\) . The union of \(U'\) (resp. \(U_1\) ) and a suitable subset of \(X - M'\) (resp. \(X - M\) ) is a fundamental domain, evidently \(\mu\) -measurable, for \(G'\) . By Integration, Chap. VII, § 2, no. 10, Cor. of Th. 4, \(\mu(U') = \mu(U_1)\) . This proves that \(\text{Card } \Lambda = (G : G')\) is finite, and that \(\mu(U') = (\text{Card } \Lambda)\mu(U)\) .

PROPOSITION 7. Assume that R is irreducible. Let \(B = \{\alpha_{1}, \ldots, \alpha_{l}\}\) be a basis of R, f the connection index of R (§ 1, no. 9) and \(\tilde{\alpha} = n_{1}\alpha_{1} + \cdots + n_{l}\alpha_{l}\) the highest root of R (for the order defined by B). Then the order of W is equal to

\[
(l!) n _ {1} n _ {2} \dots n _ {l} f.
\]

Let \((\overline{\omega}_{1},\ldots,\overline{\omega}_{l})\) be the basis of P(R \(^{\prime}\) ) dual to B. By the Cor. of Prop. 5, the open simplex C with vertices \(0,n_{1}^{-1}\overline{\omega}_{1},\ldots,n_{l}^{-1}\overline{\omega}_{l}\) is an alcove of E. Choose a Haar measure \(\mu\) on the additive group V*. Let A be the set of elements of V* of the form \(\xi_{1}\overline{\omega}_{1}+\cdots+\xi_{l}\overline{\omega}_{l}\) , with \(0<\xi_{i}<1\) for \(i=1,\ldots,l\) . By Cor. 2 of Prop. 15 of Integration, Chap. VII, § 1, no. 10,

\[
\mu (\mathrm{A}) / \mu (\mathrm{C}) = (l!) n _ {1} n _ {2} \dots n _ {l}.\tag{5}
\]

On the other hand, let \(A'\) be the set of elements of \(V^{*}\) of the form

\[
\xi_ {1} \alpha_ {1} ^ {\prime} + \dots + \xi_ {l} \alpha_ {l} ^ {\prime},
\]

with \(0 < \xi_i < 1\) for \(i = 1, \ldots, l\) . Since \((\alpha_1', \ldots, \alpha_l')\) is a basis of the \(\mathbf{Z}\) -module \(\mathrm{Q}(\mathrm{R}^{\prime})\) , we can apply Lemma 1 with \(\mathrm{X} = \mathrm{V}^*\) ; \(\mathrm{G} = \mathrm{W}_a\) , \(\mathrm{G}' = \mathrm{Q}\) , \(\mathrm{U} = \mathrm{C}\) and \(\mathrm{U}' = \mathrm{A}'\) . This gives

\[
\mu \left(\mathrm{A} ^ {\prime}\right) / \mu (\mathrm{C}) = \left(\mathrm{W} _ {a}: \mathrm{Q}\right) = \text { Card   W }.\tag{6}
\]

Finally, Lemma 1 can be applied again, taking \(\mathrm{X} = \mathrm{V}^*\) , \(\mathrm{G} = \mathrm{P}\) , \(\mathrm{G}' = \mathrm{Q}\) , \(\mathrm{U} = \mathrm{A}\) and \(\mathrm{U}' = \mathrm{A}'\) . This gives

\[
\mu (A ^ {\prime}) / \mu (A) = (P: Q) = (P (R ^ {\prime}): Q (R ^ {\prime})) = f.\tag{7}
\]

The proposition now follows by comparing formulas (5), (6) and (7).

